\begin{lstlisting}[language=Python, caption={Core algorithm for Push Gravity (PG) shell parcellation including embedded kinematic centrifugal acceleration corrections.}, label={lst:pg_core_engine}]
import numpy as np
from scipy.integrate import dblquad

def calculate_gravity_pg_core(planet_params, station_type='equator', C_coeff=0.0, omega=0.0):
    """
    Computes the effective surface gravity (g_eff) using exact shell-by-shell 
    chord parcellation under the mapping condition: (Lambda^2 * u_e0) / 4pi = G.
    Natively accounts for the kinetic centrifugal acceleration at the equator.
    
    Parameters:
    -----------
    planet_params : dict
        Contains lists for semi-major axes 'a' (m), flattenings 'f', 
        and volumetric densities 'rho' (kg/m^3), ordered from Layer 1 to 11.
    station_type : str
        'equator' or 'polar' positioning for the local coordinate origin P.
    C_coeff : float
        Dimensionless forward-scattering coefficient (C).
    omega : float
        Planetary angular rotation velocity (rad/s). Defaults to 0.0 (static).
    """
    # Universal Gravitational Constant
    G = 6.67430e-11 
    
    # Establish surface ellipsoid parameters for local origin P placement
    a_surf = planet_params['a'][-1]
    f_surf = planet_params['f'][-1]
    c_surf = a_surf * (1.0 - f_surf)
    
    if station_type == 'polar':
        X_P, Y_P, Z_P = 0.0, 0.0, c_surf
    else: # equator
        X_P, Y_P, Z_P = a_surf, 0.0, 0.0

    a = np.array(planet_params['a'])
    f = np.array(planet_params['f'])
    c = a * (1.0 - f)
    rho = np.array(planet_params['rho'])
    
    def get_chord_areal_density(alpha, beta):
        # Directional unit vectors mapped to planetary frame from station P
        if station_type == 'polar':
            ux = np.sin(alpha) * np.cos(beta)
            uy = np.sin(alpha) * np.sin(beta)
            uz = -np.cos(alpha)
        else:
            ux = -np.cos(alpha)
            uy = np.sin(alpha) * np.cos(beta)
            uz = np.sin(alpha) * np.sin(beta)
            
        s1 = np.zeros(12) 
        s2 = np.zeros(12)
        
        # Ellipsoidal intersection quadratic metric system for boundaries i=1..11
        for idx in range(11):
            i = idx + 1
            ai = a[idx]
            ci = c[idx]
            
            A_quad = (ux**2 + uy**2) / ai**2 + (uz**2) / ci**2
            B_quad = (X_P*ux + Y_P*uy) / ai**2 + (Z_P*uz) / ci**2
            C_quad = (X_P**2 + Y_P**2) / ai**2 + (Z_P**2) / ci**2 - 1.0
            
            discriminant = B_quad**2 - A_quad*C_quad
            if discriminant >= 0:
                root1 = (-B_quad - np.sqrt(discriminant)) / A_quad
                root2 = (-B_quad + np.sqrt(discriminant)) / A_quad
                s1[i] = min(root1, root2)
                s2[i] = max(root1, root2)
            else:
                s1[i] = 0.0
                s2[i] = 0.0
                
        # Enforce baseline boundary origin condition s_{1,0} = s_{2,0} = 0
        s1[0] = 0.0
        s2[0] = 0.0
        
        # Structural shell parcellation loop (Equations 1 & 2)
        total_lambda = 0.0
        for idx in range(11):
            i = idx + 1
            delta_s = s1[i-1] - s1[i] + s2[i] - s2[i-1]
            lambda_ch_i = rho[idx] * delta_s
            total_lambda += lambda_ch_i
            
        return total_lambda

    # Integral definition over the internal local horizon hemisphere
    integrand = lambda alpha, beta: get_chord_areal_density(alpha, beta) * np.cos(alpha) * np.sin(alpha)
    alpha_max = np.pi / 2.0
    
    integral_value, _ = dblquad(integrand, 0.0, 2.0*np.pi, lambda x: 0.0, lambda x: alpha_max)
    
    # Net pure gravitational acceleration mapping
    g_net = G * (1.0 - C_coeff) * integral_value
    
    # Apply kinematic centrifugal reduction at the equator if omega is supplied
    if station_type == 'equator' and omega > 0.0:
        centrifugal_acceleration = (omega**2) * a_surf
        g_eff = g_net - centrifugal_acceleration
    else:
        g_eff = g_net
        
    return g_eff
\end{lstlisting}
%
\begin{lstlisting}[language=Python, caption={Verification testbench evaluating Earth's effective surface gravity (g_eff) against true geodetic reference marks.}, label={lst:pg_verification_testbench}]
# (Requires calculate_gravity_pg_core from Listing 1)
import numpy as np

# Structured 11-Layer Earth Parameter Profiles from Layer 1 (Inner Core) to Layer 11 (Crust)
earth_prem_profile = {
    'a': [1221.5e3, 3480.0e3, 3630.0e3, 5701.0e3, 5771.0e3, 5971.0e3, 6151.0e3, 6346.0e3, 6356.0e3, 6368.0e3, 6378.1e3],
    'f': list(np.linspace(1.0/415.0, 1.0/298.257, 11)),
    'rho': [13000.0, 11500.0, 10000.0, 5200.0, 4800.0, 4400.0, 4000.0, 3600.0, 3400.0, 3000.0, 2700.0]
}

# Earth true sidereal rotation velocity metric (rad/s)
omega_earth = 7.292115e-5 

# Calibrated optimization index matching geodetic profile constraints
C_optimized = 0.037536

# Run simulation tracking
g_earth_eff_eq = calculate_gravity_pg_core(earth_prem_profile, station_type='equator', C_coeff=C_optimized, omega=omega_earth)
g_earth_eff_pol = calculate_gravity_pg_core(earth_prem_profile, station_type='polar', C_coeff=C_optimized, omega=omega_earth)

print("=== REFRESHED PUSH GRAVITY VERIFICATION OUTPUT ===")
print(f"Effective Polar Gravity (Model)   : {g_earth_eff_pol:.4f} m/s^2 (Target: 9.8322)")
print(f"Effective Equat. Gravity (Model)  : {g_earth_eff_eq:.4f} m/s^2 (Target: 9.7803)")
print(f"Effective Dynamic Delta g (Pol-Eq): {g_earth_eff_pol - g_earth_eff_eq:.4f} m/s^2")
\end{lstlisting}
%
\begin{lstlisting}[language=Python, caption={Inversion numerical protocol isolating C_opt by targeting raw, uncorrected geodetic surface data directly.}, label={lst:pg_optimizer_inversion}]
from scipy.optimize import minimize_scalar
# (Requires calculate_gravity_pg_core and earth_prem_profile from the previous listings)

def optimize_forward_scattering_direct():
    """
    Finds the optimized C coefficient by matching the geometric model outputs
    directly against Earth's raw observed effective surface gravity baselines.
    """
    # Direct observational geodetic target profiles (Raw surface measurements)
    target_g_pol = 9.8322 # Polar limit (completely centrifugal-free)
    target_g_eq  = 9.7803 # Raw equatorial surface metric
    
    omega_earth = 7.292115e-5
    
    # Pre-compute unscaled raw pure-gravitational geometric integrations at C = 0, omega = 0
    g_grav_pol_raw = calculate_gravity_pg_core(earth_prem_profile, station_type='polar', C_coeff=0.0, omega=0.0)
    g_grav_eq_raw  = calculate_gravity_pg_core(earth_prem_profile, station_type='equator', C_coeff=0.0, omega=0.0)
    
    a_surf = earth_prem_profile['a'][-1]
    centrifugal_loss = (omega_earth**2) * a_surf
    
    # Loss objective tracking through least-squares minimization
    def objective_loss(C):
        g_model_pol = (1.0 - C) * g_grav_pol_raw
        g_model_eq  = ((1.0 - C) * g_grav_eq_raw) - centrifugal_loss
        
        residual = (g_model_pol - target_g_pol)**2 + (g_model_eq - target_g_eq)**2
        return residual

    # Restrict optimization space to physically coherent bounds: 0 <= C < 1
    res = minimize_scalar(objective_loss, bounds=(0.0, 0.99), method='bounded')
    C_opt = res.x
    
    print("=== REFRESHED SCATTERING COEFFICIENT DIRECT OPTIMIZATION ===")
    print(f"Optimized Forward-Scattering C_opt : {C_opt:.6f}")
    print(f"Model Effective Equat. Gravity     : {((1.0 - C_opt) * g_grav_eq_raw) - centrifugal_loss:.4f} m/s^2")
    print(f"Model Effective Polar Gravity      : {(1.0 - C_opt) * g_grav_pol_raw:.4f} m/s^2")
    
    return C_opt

# Execute direct numerical inversion loop
C_opt_verified = optimize_forward_scattering_direct()
\end{lstlisting}
%
\lstdefinestyle{pythonstyle}{
    language=Python,
    basicstyle=\ttfamily\footnotesize,
    keywordstyle=\color{blue}\bfseries,
    commentstyle=\color{gray}\itshape,
    stringstyle=\color{teal},
    numbers=left,
    numberstyle=\tiny\color{gray},
    stepnumber=1,
    numbersep=8pt,
    backgroundcolor=\color{black!2},
    frame=single,
    rulecolor=\color{black!30},
    breaklines=true,
    breakatwhitespace=true,
    showstringspaces=false,
    tabsize=4,
    captionpos=b
}

\begin{lstlisting}[style=pythonstyle, caption={Python script for computing and plotting the scattering-corrected $q$-factor over a logarithmic opacity domain.}, label={lst:q_factor_plot}]
import numpy as np
import matplotlib.pyplot as plt

def q_factor(u, C):
    """Computes q(u, C) with a Taylor series safeguard for u -> 0."""
    u = np.asarray(u, dtype=float)
    res = np.zeros_like(u)
    
    # Use Taylor series for u < 1e-4 to avoid numerical cancellation
    small = u < 1e-4
    u_s = u[small]
    res[small] = (1.0 - C) - 0.75 * (1.0 - 2.0 * C) * u_s + 0.4 * (1.0 - 3.0 * C) * (u_s**2)
    
    # Exact analytical formula for u >= 1e-4
    u_l = u[~small]
    term1 = 1.0 / u_l
    term2 = (1.0 + 2.0 * C) / (2.0 * (u_l**3))
    term3 = np.exp(-2.0 * u_l) * (
        (2.0 * C / u_l) + ((1.0 + 2.0 * C) / (u_l**2)) + ((1.0 + 2.0 * C) / (2.0 * (u_l**3)))
    )
    res[~small] = 0.75 * (term1 - term2 + term3)
    
    return res

# Logarithmic opacity domain u in [10^-3, 10^2]
u_vals = np.logspace(-3, 2, 1000)
C_values = [0.0, 0.1, 0.2, 0.3, 0.4, 0.5]
colors = ['#000000', '#1f77b4', '#ff7f0e', '#2ca02c', '#d62728', '#9467bd']
linestyles = ['-', '--', '-.', ':', '--', '-.']

fig, ax = plt.subplots(figsize=(8, 5), dpi=300)

# Shade parameter range between C=0.0 and C=0.5
q_upper = q_factor(u_vals, 0.0)
q_lower = q_factor(u_vals, 0.5)
ax.fill_between(u_vals, q_lower, q_upper, color='steelblue', alpha=0.15,
                label=r'Range: $C \in [0, 0.5]$')

# Plot individual curves
for idx, C in enumerate(C_values):
    lw = 2.5 if C == 0.0 else 1.4
    lbl = r'$C = 0.0$ (Pure Absorption)' if C == 0.0 else f'$C = {C:.1f}$'
    ax.plot(u_vals, q_factor(u_vals, C), color=colors[idx],
            linestyle=linestyles[idx], linewidth=lw, label=lbl)

# Configure logarithmic x-axis and linear y-axis
ax.set_xscale('log')
ax.set_xlim(1e-3, 100)
ax.set_ylim(0, 1.0)
ax.set_xlabel('opacity u', fontsize=12)
ax.set_ylabel('q-factor', fontsize=12)
ax.grid(True, which='both', linestyle='--', linewidth=0.5, alpha=0.6)
ax.legend(loc='upper right', frameon=True, fontsize=10)

plt.tight_layout()

# Export vector PDF for LaTeX inclusion
plt.savefig('q_factor_plot.pdf', format='pdf', bbox_inches='tight')
plt.show()
\end{lstlisting}
%
\lstdefinestyle{pythonstyle}{
    language=Python,
    basicstyle=\ttfamily\footnotesize,
    keywordstyle=\color{blue}\bfseries,
    commentstyle=\color{gray}\itshape,
    stringstyle=\color{teal},
    numbers=left,
    numberstyle=\tiny\color{gray},
    stepnumber=1,
    numbersep=8pt,
    backgroundcolor=\color{black!2},
    frame=single,
    rulecolor=\color{black!30},
    breaklines=true,
    breakatwhitespace=true,
    showstringspaces=false,
    tabsize=4,
    captionpos=b
}

\begin{lstlisting}[style=pythonstyle, caption={Python script for plotting the scattering-corrected $q$-factor across sub-critical, critical ($C = 1.0$), and super-critical ($C > 1.0$) regimes.}, label={lst:q_factor_critical}]
import numpy as np
import matplotlib.pyplot as plt

def q_factor(u, C):
    """Computes q(u, C) with a Taylor series safeguard for u -> 0."""
    u = np.asarray(u, dtype=float)
    res = np.zeros_like(u)
    
    # Use Taylor series for u < 1e-4 to avoid numerical cancellation
    small = u < 1e-4
    u_s = u[small]
    res[small] = (1.0 - C) - 0.75 * (1.0 - 2.0 * C) * u_s + 0.4 * (1.0 - 3.0 * C) * (u_s**2)
    
    # Exact analytical formula for u >= 1e-4
    u_l = u[~small]
    term1 = 1.0 / u_l
    term2 = (1.0 + 2.0 * C) / (2.0 * (u_l**3))
    term3 = np.exp(-2.0 * u_l) * (
        (2.0 * C / u_l) + ((1.0 + 2.0 * C) / (u_l**2)) + ((1.0 + 2.0 * C) / (2.0 * (u_l**3)))
    )
    res[~small] = 0.75 * (term1 - term2 + term3)
    
    return res

# Logarithmic opacity domain u in [10^-3, 10^2]
u_vals = np.logspace(-3, 2, 1000)
C_values = [0.0, 0.25, 0.5, 0.75, 1.0, 1.2]
colors = ['#000000', '#1f77b4', '#2ca02c', '#ff7f0e', '#d62728', '#9467bd']
linestyles = ['-', '--', '-.', ':', '-', '--']

fig, ax = plt.subplots(figsize=(8.5, 5.2), dpi=300)

# Shade the strictly attractive scattering band between C=0.0 and C=1.0
q_upper = q_factor(u_vals, 0.0)
q_crit = q_factor(u_vals, 1.0)
ax.fill_between(u_vals, q_crit, q_upper, color='steelblue', alpha=0.12,
                label=r'Attractive Band: $C \in [0, 1.0]$')

# Plot individual curves
for idx, C in enumerate(C_values):
    lw = 2.5 if C in (0.0, 1.0) else 1.4
    if C == 0.0:
        lbl = r'$C = 0.00$ (Pure Absorption)'
    elif C == 1.0:
        lbl = r'$C = 1.00$ (Critical Cancellation)'
    elif C == 1.2:
        lbl = r'$C = 1.20$ (Repulsive Thin Limit)'
    else:
        lbl = f'$C = {C:.2f}$'
        
    ax.plot(u_vals, q_factor(u_vals, C), color=colors[idx],
            linestyle=linestyles[idx], linewidth=lw, label=lbl)

# Horizontal zero-reference line separating attractive (q > 0) and repulsive (q < 0) regimes
ax.axhline(0.0, color='black', linestyle=':', linewidth=1.0, alpha=0.7)

# Configure logarithmic x-axis and linear y-axis
ax.set_xscale('log')
ax.set_xlim(1e-3, 100)
ax.set_ylim(-0.3, 1.0)
ax.set_xlabel('opacity u', fontsize=12)
ax.set_ylabel('q-factor', fontsize=12)
ax.grid(True, which='both', linestyle='--', linewidth=0.5, alpha=0.6)
ax.legend(loc='upper right', frameon=True, fontsize=9.5)

plt.tight_layout()

# Export vector PDF for LaTeX inclusion
plt.savefig('q_factor_critical_regimes.pdf', format='pdf', bbox_inches='tight')
plt.show()
\end{lstlisting}
%
\begin{lstlisting}[language=Python, caption={Complete Push Gravity (PG) continuous integration engine using original PREM polynomials (Dziewonski \& Anderson, 1981) and branch-free exact shell segmentation (C = 0).}, label={lst:pg_prem_continuous_complete}]
import numpy as np
from scipy.integrate import dblquad

# Original 11-Layer PREM Polynomial Coefficients (Dziewonski & Anderson, 1981)
# Ordered from Layer 1 (Inner Core) to Layer 11 (Upper Crust / Ocean avg)
# rho_i(x) = c0 + c1*x + c2*x^2 + c3*x^3 [kg/m^3], where x = r_eq / a_surf
prem_continuous_11 = {
    'a': [1221.5e3, 3480.0e3, 3630.0e3, 5701.0e3, 5771.0e3, 
          5971.0e3, 6151.0e3, 6291.0e3, 6346.6e3, 6356.0e3, 6371.0e3],
    'f': list(np.linspace(1.0/415.0, 1.0/298.257, 11)),
    'poly': [
        [13088.5,     0.0,  -8838.1,     0.0], # 1: Inner Core (0 - 1221.5 km)
        [12581.5, -1263.8,  -3642.6, -5528.1], # 2: Outer Core (1221.5 - 3480 km)
        [ 7956.5, -6476.1,   5528.3, -3080.7], # 3: D'' Layer (3480 - 3630 km)
        [ 7956.5, -6476.1,   5528.3, -3080.7], # 4: Lower Mantle Base (3630 - 5701 km)
        [ 5319.7, -1483.6,      0.0,     0.0], # 5: Lower Mantle Upper (5701 - 5771 km)
        [11249.4, -8029.8,      0.0,     0.0], # 6: Transition Zone (5771 - 5971 km)
        [ 7108.9, -3804.5,      0.0,     0.0], # 7: Transition Zone (5971 - 6151 km)
        [ 2691.0,   692.4,      0.0,     0.0], # 8: LVZ / LID (6151 - 6291 km)
        [ 2691.0,   692.4,      0.0,     0.0], # 9: Upper Mantle (6291 - 6346.6 km)
        [ 2900.0,     0.0,      0.0,     0.0], # 10: Lower Crust (6346.6 - 6356 km)
        [ 2600.0,     0.0,      0.0,     0.0]  # 11: Upper Crust / Ocean (6356 - 6371 km)
    ]
}

# 5-point Gauss-Legendre quadrature nodes and weights (exact up to degree 9)
GL_NODES, GL_WEIGHTS = np.polynomial.legendre.leggauss(5)

def calculate_gravity_pg_prem_continuous(planet_params, station_type='equator', omega=0.0):
    """
    Computes effective surface gravity (g_eff) with C = 0 using continuous space 
    integration along ellipsoidal chords and original PREM density polynomials.
    Preserves the exact branch-free boundary segmentation algorithm.
    """
    G = 6.67430e-11
    a = np.array(planet_params['a'])
    f = np.array(planet_params['f'])
    c = a * (1.0 - f)
    polys = np.array(planet_params['poly'])
    
    a_surf = a[-1]
    c_surf = c[-1]
    
    # Station P coordinates in global Cartesian frame
    if station_type == 'polar':
        X_P, Y_P, Z_P = 0.0, 0.0, c_surf
    else:
        X_P, Y_P, Z_P = a_surf, 0.0, 0.0

    def eval_rho_layer(idx, s_val, ux, uy, uz):
        # 3D position vector X(s) = X_P + s * u
        X = X_P + s_val * ux
        Y = Y_P + s_val * uy
        Z = Z_P + s_val * uz
        # Normalized equivalent ellipsoidal radius x (unclamped for exact algebraic cancellation)
        fi = f[idx]
        r_eq = np.sqrt(X**2 + Y**2 + (Z / (1.0 - fi))**2)
        x = r_eq / a_surf
        c0, c1, c2, c3 = polys[idx]
        return c0 + c1*x + c2*(x**2) + c3*(x**3)

    def integrate_oriented_segment(idx, s_start, s_end, ux, uy, uz):
        # Oriented line integral from s_start to s_end preserving sign
        half_len = 0.5 * (s_end - s_start)
        if half_len == 0.0:
            return 0.0
        mid_pt = 0.5 * (s_end + s_start)
        acc = 0.0
        for node, weight in zip(GL_NODES, GL_WEIGHTS):
            s_q = mid_pt + half_len * node
            acc += weight * eval_rho_layer(idx, s_q, ux, uy, uz)
        return half_len * acc

    def get_continuous_areal_density(alpha, beta):
        # Directional unit vector mapped from local normal at P
        if station_type == 'polar':
            ux = np.sin(alpha) * np.cos(beta)
            uy = np.sin(alpha) * np.sin(beta)
            uz = -np.cos(alpha)
        else:
            ux = -np.cos(alpha)
            uy = np.sin(alpha) * np.cos(beta)
            uz = np.sin(alpha) * np.sin(beta)
            
        s1 = np.zeros(12)
        s2 = np.zeros(12)
        
        # Solve quadratic intersection A_i*s^2 + 2*B_i*s + C_i = 0 for i = 1..11
        for idx in range(11):
            i = idx + 1
            ai, ci = a[idx], c[idx]
            A_q = (ux**2 + uy**2)/ai**2 + (uz**2)/ci**2
            B_q = (X_P*ux + Y_P*uy)/ai**2 + (Z_P*uz)/ci**2
            C_q = (X_P**2 + Y_P**2)/ai**2 + (Z_P**2)/ci**2 - 1.0
            disc = B_q**2 - A_q*C_q
            if disc >= 0.0:
                sqrt_d = np.sqrt(disc)
                s1[i] = (-B_q - sqrt_d) / A_q
                s2[i] = (-B_q + sqrt_d) / A_q
                
        # Exact continuous counterpart of Equations (1) and (2) without if/else branches
        total_lambda = 0.0
        for idx in range(11):
            i = idx + 1
            lambda_ch_i = (
                integrate_oriented_segment(idx, s1[i], s1[i-1], ux, uy, uz) +
                integrate_oriented_segment(idx, s2[i-1], s2[i], ux, uy, uz)
            )
            total_lambda += lambda_ch_i
            
        return total_lambda

    # Double angular quadrature over the visible internal hemisphere
    integrand = lambda alpha, beta: get_continuous_areal_density(alpha, beta) * np.cos(alpha) * np.sin(alpha)
    integral_val, _ = dblquad(
        integrand, 0.0, 2.0*np.pi, lambda x: 0.0, lambda x: np.pi/2.0,
        epsabs=1e-5, epsrel=1e-5
    )
    
    # Unscattered baseline (C = 0 -> 1 - C = 1)
    g_net = G * integral_val
    
    # Kinematic centrifugal subtraction at equator
    if station_type == 'equator' and omega > 0.0:
        return g_net - (omega**2) * a_surf
    return g_net

# --- EXECUTION TESTBENCH (C = 0) ---
omega_earth = 7.292115e-5
g_pol_c0 = calculate_gravity_pg_prem_continuous(prem_continuous_11, station_type='polar', omega=omega_earth)
g_eq_c0  = calculate_gravity_pg_prem_continuous(prem_continuous_11, station_type='equator', omega=omega_earth)

print(f"Continuous PREM (C=0) Polar Gravity      : {g_pol_c0:.4f} m/s^2")
print(f"Continuous PREM (C=0) Equatorial Gravity : {g_eq_c0:.4f} m/s^2")
\end{lstlisting}
%
\begin{lstlisting}[language=Python, caption={Generalized N-layer Push Gravity (PG) integration script with automated bulk mass and moment-of-inertia verification (C = 0).}, label={lst:pg_inner_planets_verified}]
import numpy as np
from scipy.integrate import dblquad

def get_flattening_profile(planet_params):
    """Returns the array of layer flattenings f_i according to Clairaut scaling."""
    n_layers = len(planet_params['a'])
    if 'f' in planet_params:
        return np.array(planet_params['f'], dtype=float)
    f_surf = float(planet_params['f_surf'])
    return np.linspace(0.72 * f_surf, f_surf, n_layers) if f_surf > 0.0 else np.zeros(n_layers)

def verify_planetary_constraints(planet_params):
    """
    Computes bulk mass M and normalized polar moment of inertia I/(M*R_eq^2)
    for nested oblate spheroidal shells to verify geodetic consistency.
    """
    a = np.array(planet_params['a'], dtype=float)
    rho = np.array(planet_params['rho'], dtype=float)
    f = get_flattening_profile(planet_params)
    c = a * (1.0 - f)
    
    total_mass = 0.0
    total_inertia_z = 0.0
    
    for idx in range(len(a)):
        # Outer boundary volume and polar moment of inertia factor of ellipsoid
        V_out = (4.0 / 3.0) * np.pi * (a[idx]**2) * c[idx]
        Iz_out = 0.4 * V_out * (a[idx]**2)
        
        if idx == 0:
            V_in, Iz_in = 0.0, 0.0
        else:
            V_in = (4.0 / 3.0) * np.pi * (a[idx-1]**2) * c[idx-1]
            Iz_in = 0.4 * V_in * (a[idx-1]**2)
            
        shell_mass = rho[idx] * (V_out - V_in)
        shell_inertia = rho[idx] * (Iz_out - Iz_in)
        
        total_mass += shell_mass
        total_inertia_z += shell_inertia
        
    moi_factor = total_inertia_z / (total_mass * (a[-1]**2))
    return total_mass, moi_factor

def calculate_gravity_pg_generalized(planet_params, station_type='equator', C_coeff=0.0):
    """
    Computes effective surface gravity (g_eff) using the exact shell-by-shell 
    chord segmentation algorithm for an arbitrary N-layer oblate spheroid.
    """
    G = 6.67430e-11
    a = np.array(planet_params['a'], dtype=float)
    rho = np.array(planet_params['rho'], dtype=float)
    f = get_flattening_profile(planet_params)
    c = a * (1.0 - f)
    
    omega = float(planet_params['omega'])
    n_layers = len(a)
    a_surf, c_surf = a[-1], c[-1]
    
    # Origin station P coordinates in global planetary frame
    if station_type == 'polar':
        X_P, Y_P, Z_P = 0.0, 0.0, c_surf
    else:
        X_P, Y_P, Z_P = a_surf, 0.0, 0.0

    def get_chord_areal_density(alpha, beta):
        # Directional unit vectors relative to local internal normal at P
        if station_type == 'polar':
            ux = np.sin(alpha) * np.cos(beta)
            uy = np.sin(alpha) * np.sin(beta)
            uz = -np.cos(alpha)
        else:
            ux = -np.cos(alpha)
            uy = np.sin(alpha) * np.cos(beta)
            uz = np.sin(alpha) * np.sin(beta)
            
        s1 = np.zeros(n_layers + 1)
        s2 = np.zeros(n_layers + 1)
        
        # Solve quadratic boundary intersections for layers i = 1..N
        for idx in range(n_layers):
            i = idx + 1
            ai, ci = a[idx], c[idx]
            
            A_q = (ux**2 + uy**2) / ai**2 + (uz**2) / ci**2
            B_q = (X_P*ux + Y_P*uy) / ai**2 + (Z_P*uz) / ci**2
            C_q = (X_P**2 + Y_P**2) / ai**2 + (Z_P**2) / ci**2 - 1.0
            
            disc = B_q**2 - A_q*C_q
            if disc >= 0.0:
                sqrt_d = np.sqrt(disc)
                r1 = (-B_q - sqrt_d) / A_q
                r2 = (-B_q + sqrt_d) / A_q
                s1[i] = min(r1, r2)
                s2[i] = max(r1, r2)
                
        # Exact shell-by-shell segmentation (Equations 1 & 2)
        total_lambda = 0.0
        for idx in range(n_layers):
            i = idx + 1
            delta_s = s1[i-1] - s1[i] + s2[i] - s2[i-1]
            total_lambda += rho[idx] * delta_s
            
        return total_lambda

    # Double angular integration over the visible internal hemisphere
    integrand = lambda alpha, beta: get_chord_areal_density(alpha, beta) * np.cos(alpha) * np.sin(alpha)
    integral_val, _ = dblquad(
        integrand, 0.0, 2.0*np.pi, lambda x: 0.0, lambda x: np.pi/2.0,
        epsabs=1e-5, epsrel=1e-5
    )
    
    # Pure gravitational acceleration under direct mapping (Lambda^2 * u_e0 / 4pi = G)
    g_net = G * (1.0 - C_coeff) * integral_val
    
    # Kinematic centrifugal subtraction at the equator
    if station_type == 'equator' and omega != 0.0:
        g_eff = g_net - (omega**2) * a_surf
    else:
        g_eff = g_net
        
    return g_eff

# --- MASS- AND INERTIA-CONSTRAINED INNER SOLAR SYSTEM DATASETS ---
inner_planets_data = {
    'Moon': {
        'a': [240.0e3, 380.0e3, 1693.1e3, 1738.1e3],
        'rho': [7398.0, 5194.0, 3379.0, 2595.0],
        'f_surf': 0.0012,
        'omega': 2.662e-6
    },
    'Mercury': {
        'a': [2020.0e3, 2399.7e3, 2439.7e3],
        'rho': [7197.0, 3159.0, 2755.0],
        'f_surf': 0.0009,
        'omega': 1.240e-6
    },
    'Venus': {
        'a': [3200.0e3, 5390.0e3, 6021.8e3, 6051.8e3],
        'rho': [10827.0, 4672.0, 3550.0, 2900.0],
        'f_surf': 0.0000,
        'omega': -2.992e-7
    },
    'Mars': {
        'a': [600.0e3, 1680.0e3, 1835.0e3, 3354.2e3, 3396.2e3],
        'rho': [7500.0, 6500.0, 4400.0, 3575.0, 2700.0],
        'f_surf': 0.00589,
        'omega': 7.088e-5
    }
}

# --- EXECUTION & GEODETIC CONSTRAINT VERIFICATION LOOP (C = 0) ---
print("=== PUSH GRAVITY: INNER SOLAR SYSTEM VERIFICATION (C = 0) ===")
for name, params in inner_planets_data.items():
    mass, moi = verify_planetary_constraints(params)
    g_eq  = calculate_gravity_pg_generalized(params, station_type='equator', C_coeff=0.0)
    g_pol = calculate_gravity_pg_generalized(params, station_type='polar', C_coeff=0.0)
    print(f"{name:8s} | M: {mass:.4e} kg | I/MR^2: {moi:.4f} | g_eq: {g_eq:.3f} | g_pol: {g_pol:.3f} m/s^2")
\end{lstlisting}
%
\begin{lstlisting}[language=Python, caption={Generalized N-layer Push Gravity (PG) integration script with automated bulk mass, moment-of-inertia, and Newtonian J2 multipole benchmarking (C = 0).}, label={lst:pg_inner_planets_full}]
import numpy as np
from scipy.integrate import dblquad

def get_flattening_profile(planet_params):
    """Returns the array of layer flattenings f_i according to Clairaut scaling."""
    n_layers = len(planet_params['a'])
    if 'f' in planet_params:
        return np.array(planet_params['f'], dtype=float)
    f_surf = float(planet_params['f_surf'])
    return np.linspace(0.72 * f_surf, f_surf, n_layers) if f_surf > 0.0 else np.zeros(n_layers)

def verify_planetary_constraints(planet_params):
    """
    Computes bulk mass M, normalized polar moment of inertia I_z/(M*R_eq^2),
    and the dynamical zonal harmonic J2 for nested oblate spheroidal shells.
    """
    a = np.array(planet_params['a'], dtype=float)
    rho = np.array(planet_params['rho'], dtype=float)
    f = get_flattening_profile(planet_params)
    c = a * (1.0 - f)
    
    total_mass = 0.0
    total_Iz = 0.0
    total_Ix = 0.0
    
    for idx in range(len(a)):
        V_out = (4.0 / 3.0) * np.pi * (a[idx]**2) * c[idx]
        Iz_out = 0.4 * V_out * (a[idx]**2)
        Ix_out = 0.2 * V_out * (a[idx]**2 + c[idx]**2)
        
        if idx == 0:
            V_in, Iz_in, Ix_in = 0.0, 0.0, 0.0
        else:
            V_in = (4.0 / 3.0) * np.pi * (a[idx-1]**2) * c[idx-1]
            Iz_in = 0.4 * V_in * (a[idx-1]**2)
            Ix_in = 0.2 * V_in * (a[idx-1]**2 + c[idx-1]**2)
            
        total_mass += rho[idx] * (V_out - V_in)
        total_Iz += rho[idx] * (Iz_out - Iz_in)
        total_Ix += rho[idx] * (Ix_out - Ix_in)
        
    a_surf = a[-1]
    moi_factor = total_Iz / (total_mass * (a_surf**2))
    # MacCullagh's formula for zonal harmonic J2
    J2 = (total_Iz - total_Ix) / (total_mass * (a_surf**2))
    
    return total_mass, moi_factor, J2

def compute_newtonian_j2_gravity(planet_params, mass, J2):
    """
    Computes standard macroscopic Newtonian surface gravity at the equator 
    and poles using first-order spherical harmonic expansion up to J2.
    """
    G = 6.67430e-11
    a_surf = float(planet_params['a'][-1])
    f_surf = float(planet_params['f_surf'])
    c_surf = a_surf * (1.0 - f_surf)
    omega = float(planet_params['omega'])
    
    # Equatorial effective gravity (including centrifugal reduction)
    g_eq_grav = (G * mass / (a_surf**2)) * (1.0 + 1.5 * J2)
    g_eq_newt = g_eq_grav - (omega**2) * a_surf
    
    # Polar gravity
    g_pol_newt = (G * mass / (c_surf**2)) * (1.0 - 3.0 * J2)
    
    return g_eq_newt, g_pol_newt

def calculate_gravity_pg_generalized(planet_params, station_type='equator', C_coeff=0.0):
    """
    Computes effective surface gravity (g_eff) using the exact shell-by-shell 
    chord segmentation algorithm for an arbitrary N-layer oblate spheroid.
    """
    G = 6.67430e-11
    a = np.array(planet_params['a'], dtype=float)
    rho = np.array(planet_params['rho'], dtype=float)
    f = get_flattening_profile(planet_params)
    c = a * (1.0 - f)
    
    omega = float(planet_params['omega'])
    n_layers = len(a)
    a_surf, c_surf = a[-1], c[-1]
    
    if station_type == 'polar':
        X_P, Y_P, Z_P = 0.0, 0.0, c_surf
    else:
        X_P, Y_P, Z_P = a_surf, 0.0, 0.0

    def get_chord_areal_density(alpha, beta):
        if station_type == 'polar':
            ux = np.sin(alpha) * np.cos(beta)
            uy = np.sin(alpha) * np.sin(beta)
            uz = -np.cos(alpha)
        else:
            ux = -np.cos(alpha)
            uy = np.sin(alpha) * np.cos(beta)
            uz = np.sin(alpha) * np.sin(beta)
            
        s1 = np.zeros(n_layers + 1)
        s2 = np.zeros(n_layers + 1)
        
        for idx in range(n_layers):
            i = idx + 1
            ai, ci = a[idx], c[idx]
            
            A_q = (ux**2 + uy**2) / ai**2 + (uz**2) / ci**2
            B_q = (X_P*ux + Y_P*uy) / ai**2 + (Z_P*uz) / ci**2
            C_q = (X_P**2 + Y_P**2) / ai**2 + (Z_P**2) / ci**2 - 1.0
            
            disc = B_q**2 - A_q*C_q
            if disc >= 0.0:
                sqrt_d = np.sqrt(disc)
                r1 = (-B_q - sqrt_d) / A_q
                r2 = (-B_q + sqrt_d) / A_q
                s1[i] = min(r1, r2)
                s2[i] = max(r1, r2)
                
        total_lambda = 0.0
        for idx in range(n_layers):
            i = idx + 1
            delta_s = s1[i-1] - s1[i] + s2[i] - s2[i-1]
            total_lambda += rho[idx] * delta_s
            
        return total_lambda

    integrand = lambda alpha, beta: get_chord_areal_density(alpha, beta) * np.cos(alpha) * np.sin(alpha)
    integral_val, _ = dblquad(
        integrand, 0.0, 2.0*np.pi, lambda x: 0.0, lambda x: np.pi/2.0,
        epsabs=1e-5, epsrel=1e-5
    )
    
    g_net = G * (1.0 - C_coeff) * integral_val
    if station_type == 'equator' and omega != 0.0:
        return g_net - (omega**2) * a_surf
    return g_net

# --- MASS- AND INERTIA-CONSTRAINED INNER SOLAR SYSTEM DATASETS ---
inner_planets_data = {
    'Moon': {
        'a': [240.0e3, 380.0e3, 1693.1e3, 1738.1e3],
        'rho': [7398.0, 5194.0, 3379.0, 2595.0],
        'f_surf': 0.0012,
        'omega': 2.662e-6
    },
    'Mercury': {
        'a': [2020.0e3, 2399.7e3, 2439.7e3],
        'rho': [7197.0, 3159.0, 2755.0],
        'f_surf': 0.0009,
        'omega': 1.240e-6
    },
    'Venus': {
        'a': [3200.0e3, 5390.0e3, 6021.8e3, 6051.8e3],
        'rho': [10827.0, 4672.0, 3550.0, 2900.0],
        'f_surf': 0.0000,
        'omega': -2.992e-7
    },
    'Mars': {
        'a': [600.0e3, 1680.0e3, 1835.0e3, 3354.2e3, 3396.2e3],
        'rho': [7500.0, 6500.0, 4400.0, 3575.0, 2700.0],
        'f_surf': 0.00589,
        'omega': 7.088e-5
    }
}

# --- EXECUTION & BENCHMARKING LOOP (C = 0) ---
print("=== PUSH GRAVITY VS. NEWTONIAN J2 BENCHMARK (C = 0) ===")
for name, params in inner_planets_data.items():
    mass, moi, J2 = verify_planetary_constraints(params)
    g_eq_pg   = calculate_gravity_pg_generalized(params, station_type='equator', C_coeff=0.0)
    g_pol_pg  = calculate_gravity_pg_generalized(params, station_type='polar', C_coeff=0.0)
    g_eq_newt, g_pol_newt = compute_newtonian_j2_gravity(params, mass, J2)
    
    print(f"{name:8s} | M: {mass:.3e} | I/MR^2: {moi:.4f} | J2: {J2:.2e}")
    print(f"         -> PG   : g_eq = {g_eq_pg:.3f} m/s^2, g_pol = {g_pol_pg:.3f} m/s^2")
    print(f"         -> Newt : g_eq = {g_eq_newt:.3f} m/s^2, g_pol = {g_pol_newt:.3f} m/s^2")
\end{lstlisting}
%
\begin{lstlisting}[language=Python, caption={Continuous space Push Gravity (PG) integration engine for Earth using original 11-layer PREM cubic density polynomials (Dziewonski \& Anderson, 1981) in the unscattered limit (C = 0).}, label={lst:pg_earth_prem_continuous}]
import numpy as np
from scipy.integrate import dblquad

# Original 11-Layer PREM Polynomial Coefficients (Dziewonski & Anderson, 1981)
# Ordered from Layer 1 (Inner Core) to Layer 11 (Upper Crust / Ocean avg)
# rho_i(x) = c0 + c1*x + c2*x^2 + c3*x^3 [kg/m^3], where x = r_eq / a_surf
prem_continuous_11 = {
    'a': [1221.5e3, 3480.0e3, 3630.0e3, 5701.0e3, 5771.0e3, 
          5971.0e3, 6151.0e3, 6291.0e3, 6346.6e3, 6356.0e3, 6371.0e3],
    'f': list(np.linspace(1.0/415.0, 1.0/298.257, 11)),
    'poly': [
        [13088.5,     0.0,  -8838.1,     0.0], # 1: Inner Core (0 - 1221.5 km)
        [12581.5, -1263.8,  -3642.6, -5528.1], # 2: Outer Core (1221.5 - 3480 km)
        [ 7956.5, -6476.1,   5528.3, -3080.7], # 3: D'' Layer (3480 - 3630 km)
        [ 7956.5, -6476.1,   5528.3, -3080.7], # 4: Lower Mantle Base (3630 - 5701 km)
        [ 5319.7, -1483.6,      0.0,     0.0], # 5: Lower Mantle Upper (5701 - 5771 km)
        [11249.4, -8029.8,      0.0,     0.0], # 6: Transition Zone (5771 - 5971 km)
        [ 7108.9, -3804.5,      0.0,     0.0], # 7: Transition Zone (5971 - 6151 km)
        [ 2691.0,   692.4,      0.0,     0.0], # 8: LVZ / LID (6151 - 6291 km)
        [ 2691.0,   692.4,      0.0,     0.0], # 9: Upper Mantle (6291 - 6346.6 km)
        [ 2900.0,     0.0,      0.0,     0.0], # 10: Lower Crust (6346.6 - 6356 km)
        [ 2600.0,     0.0,      0.0,     0.0]  # 11: Upper Crust / Ocean (6356 - 6371 km)
    ]
}

# 5-point Gauss-Legendre quadrature nodes and weights (exact up to degree 9)
GL_NODES, GL_WEIGHTS = np.polynomial.legendre.leggauss(5)

def calculate_gravity_pg_prem_continuous(planet_params, station_type='equator', omega=0.0):
    """
    Computes effective surface gravity (g_eff) with C = 0 using continuous space 
    integration along ellipsoidal chords and original PREM density polynomials.
    Preserves the exact branch-free boundary segmentation algorithm.
    """
    G = 6.67430e-11
    a = np.array(planet_params['a'], dtype=float)
    f = np.array(planet_params['f'], dtype=float)
    c = a * (1.0 - f)
    polys = np.array(planet_params['poly'], dtype=float)
    
    a_surf = a[-1]
    c_surf = c[-1]
    
    # Station P coordinates in global Cartesian frame
    if station_type == 'polar':
        X_P, Y_P, Z_P = 0.0, 0.0, c_surf
    else:
        X_P, Y_P, Z_P = a_surf, 0.0, 0.0

    def eval_rho_layer(idx, s_val, ux, uy, uz):
        # 3D position vector X(s) = X_P + s * u
        X = X_P + s_val * ux
        Y = Y_P + s_val * uy
        Z = Z_P + s_val * uz
        # Normalized equivalent ellipsoidal radius x (unclamped for exact algebraic cancellation)
        fi = f[idx]
        r_eq = np.sqrt(X**2 + Y**2 + (Z / (1.0 - fi))**2)
        x = r_eq / a_surf
        c0, c1, c2, c3 = polys[idx]
        return c0 + c1*x + c2*(x**2) + c3*(x**3)

    def integrate_oriented_segment(idx, s_start, s_end, ux, uy, uz):
        # Oriented line integral from s_start to s_end preserving sign
        half_len = 0.5 * (s_end - s_start)
        if half_len == 0.0:
            return 0.0
        mid_pt = 0.5 * (s_end + s_start)
        acc = 0.0
        for node, weight in zip(GL_NODES, GL_WEIGHTS):
            s_q = mid_pt + half_len * node
            acc += weight * eval_rho_layer(idx, s_q, ux, uy, uz)
        return half_len * acc

    def get_continuous_areal_density(alpha, beta):
        # Directional unit vector mapped from local normal at P
        if station_type == 'polar':
            ux = np.sin(alpha) * np.cos(beta)
            uy = np.sin(alpha) * np.sin(beta)
            uz = -np.cos(alpha)
        else:
            ux = -np.cos(alpha)
            uy = np.sin(alpha) * np.cos(beta)
            uz = np.sin(alpha) * np.sin(beta)
            
        s1 = np.zeros(12)
        s2 = np.zeros(12)
        
        # Solve quadratic intersection A_i*s^2 + 2*B_i*s + C_i = 0 for i = 1..11
        for idx in range(11):
            i = idx + 1
            ai, ci = a[idx], c[idx]
            A_q = (ux**2 + uy**2)/ai**2 + (uz**2)/ci**2
            B_q = (X_P*ux + Y_P*uy)/ai**2 + (Z_P*uz)/ci**2
            C_q = (X_P**2 + Y_P**2)/ai**2 + (Z_P**2)/ci**2 - 1.0
            disc = B_q**2 - A_q*C_q
            if disc >= 0.0:
                sqrt_d = np.sqrt(disc)
                s1[i] = (-B_q - sqrt_d) / A_q
                s2[i] = (-B_q + sqrt_d) / A_q
                
        # Exact continuous counterpart of Equations (1) and (2) without if/else branches
        total_lambda = 0.0
        for idx in range(11):
            i = idx + 1
            lambda_ch_i = (
                integrate_oriented_segment(idx, s1[i], s1[i-1], ux, uy, uz) +
                integrate_oriented_segment(idx, s2[i-1], s2[i], ux, uy, uz)
            )
            total_lambda += lambda_ch_i
            
        return total_lambda

    # Double angular quadrature over the visible internal hemisphere
    integrand = lambda alpha, beta: get_continuous_areal_density(alpha, beta) * np.cos(alpha) * np.sin(alpha)
    integral_val, _ = dblquad(
        integrand, 0.0, 2.0*np.pi, lambda x: 0.0, lambda x: np.pi/2.0,
        epsabs=1e-5, epsrel=1e-5
    )
    
    # Unscattered baseline (C = 0 -> 1 - C = 1)
    g_net = G * integral_val
    
    # Kinematic centrifugal subtraction at equator
    if station_type == 'equator' and omega > 0.0:
        return g_net - (omega**2) * a_surf
    return g_net

# --- EXECUTION TESTBENCH (C = 0) ---
omega_earth = 7.292115e-5
g_pol_c0 = calculate_gravity_pg_prem_continuous(prem_continuous_11, station_type='polar', omega=omega_earth)
g_eq_c0  = calculate_gravity_pg_prem_continuous(prem_continuous_11, station_type='equator', omega=omega_earth)

print("=== EARTH CONTINUOUS POLYNOMIAL PREM (C = 0) ===")
print(f"Effective Polar Gravity      : {g_pol_c0:.4f} m/s^2 (Ref: 9.8322)")
print(f"Effective Equatorial Gravity : {g_eq_c0:.4f} m/s^2 (Ref: 9.7803)")
\end{lstlisting}

